\chapter{Introdução}
\label{chap_introducao}

O transporte marítimo responde por mais de 80\% do comércio global de mercadorias, configurando-se como o meio mais eficiente para o transporte de cargas em grande escala \citep{unctad_transport_2025}. No entanto, o transporte marítimo internacional também contribui de forma relevante para os problemas ambientais globais, uma vez que os combustíveis utilizados nos navios mercantes são predominantemente \emph{bunker fuel}, combustível fóssil derivado de hidrocarbonetos que libera gases nocivos na atmosfera, como óxidos de enxofre ($SO_x$) e dióxido de carbono ($CO_2$) \citep{international_maritime_organization_fourth_2020,van_global_2019}.

Em razão da importância estratégica que exerce nas relações comerciais e, portanto, na economia mundial, o transporte marítimo integra os debates sobre redução de emissões de \gls{GEE} promovidos por organismos multilaterais, entre eles a \gls{IMO}, responsável pela regulação do setor \citep{international_maritime_organization_fourth_2020,imo2023strategy}. O setor enfrenta, assim, desafios significativos para atingir as metas de redução da intensidade de carbono, especialmente diante das restrições impostas pela \gls{IMO} 2020. Essas metas estabelecem objetivos para 2030 e 2050 \citep{imo_imo_2020,imo2023strategy}, conforme descrito a seguir:

\begin{itemize}
  \item Metas para 2030: reduzir a intensidade de carbono em 40\% em relação aos níveis de 2008, promovendo maior eficiência energética e o uso de combustíveis de baixa emissão.
  \item Metas para 2050: reduzir as emissões absolutas em 50\% em comparação a 2008 e alcançar a neutralidade de carbono por meio de combustíveis de emissão zero e tecnologias limpas.
\end{itemize}

As opções para o atingimento do objetivo de descarbonização incluem intervenções tecnológicas e inovações voltadas à melhoria da eficiência energética e operacional dos meios marítimos. Contudo, para alcançar as metas definidas pela \gls{IMO}, torna-se necessária a substituição dos combustíveis fósseis por fontes alternativas de energia \citep{imo2023strategy,international_council_on_clean_transportation_potential_2021}.

O hidrogênio verde é apontado na literatura como referência relevante de longo prazo para o atingimento das metas de descarbonização do transporte marítimo até 2050, por apresentar potencial de emissão zero de carbono quando produzido a partir de fontes renováveis. Quando utilizado em células de combustível, elimina as emissões de $CO_2$, enxofre ($SO_x$) e particulados. No entanto, sua adoção enfrenta desafios no curto prazo, tais como:

\begin{itemize}
  \item Infraestrutura: requer o desenvolvimento de sistemas para produção, armazenamento e abastecimento, envolvendo investimentos elevados e prazos extensos \citep{lin_alternative_2023};
  \item Custo: apresenta custo elevado atualmente, em razão do alto custo de produção por eletrólise \citep{hansson_potential_2020};
  \item Tecnologia dos navios: exige novos motores ou células de combustível, além de tanques especializados para armazenamento, o que dificulta sua implementação em larga escala até 2030 \citep{tornatore_ammonia_2022}.
\end{itemize}


Enquanto essa alternativa não for viável, o setor marítimo precisa reduzir emissões de \gls{GEE} em consonância com metas internacionais, mantendo parâmetros operacionais e custos compatíveis com sua função no comércio global. Nesse intervalo, o biobunker, definido como a mistura de biodiesel (biocombustível produzido a partir de fontes renováveis, como óleos vegetais e gorduras animais, que apresenta caráter biodegradável \citep{paulauskiene_alternative_2019}) com \emph{bunker fuel}, surge na literatura como combustível de transição com potencial para reduzir parcialmente as emissões antes da adoção mais ampla de tecnologias de emissão zero \citep{irena_pathway_2021,roux_review_2024}.

Em contextos específicos, essa alternativa pode reduzir emissões de $CO_2$, $SO_x$ e particulados, ser utilizada em motores atuais com pouca ou nenhuma modificação e, quando produzida a partir de resíduos, contribuir para a economia circular e para a redução de impactos ambientais \citep{bilgili_comparative_2021,altarazi_effects_2022,emaf_towards_2013}.

Embora não seja capaz de atingir a neutralidade de carbono proposta para 2050, o biobunker pode contribuir para a redução de emissões no curto prazo em razão da infraestrutura existente, da relativa facilidade de adoção e, em alguns contextos, da menor barreira de implementação em comparação com rotas mais disruptivas \citep{international_council_on_clean_transportation_potential_2021,altarazi_effects_2022,petzold_operation_2011}. Além disso, pode representar alternativa mais acessível para o atendimento das metas intermediárias da \gls{IMO} até 2030, especialmente em países com abundância de biomassa, como o Brasil \citep{carvalho_biofuels_2021}.

Embora as metas de descarbonização do transporte marítimo sejam definidas em âmbito global pela \gls{IMO}, a adoção de combustíveis alternativos no curto prazo é condicionada por instrumentos regulatórios regionais e por exigências de mercados específicos. Em particular, rotas internacionais que conectam países produtores de combustíveis, como o Brasil, a grandes mercados consumidores, como a União Europeia, estão sujeitas a pressões regulatórias externas que influenciam as decisões operacionais dos armadores, independentemente da bandeira do navio ou do local de abastecimento do combustível \citep{eu_reg_2023_957,eu_reg_2023_1805}.

Nesse contexto, a análise do biobunker como combustível de transição no transporte marítimo internacional operando a partir do Brasil adquire relevância não apenas por suas características técnicas e ambientais, mas também por sua potencial adequação a exigências regulatórias internacionais que afetam a competitividade das rotas marítimas \citep{carvalho_biofuels_2021,eu_reg_2023_957}.

Por essa razão, este estudo adota como referência um cenário de navegação internacional entre portos brasileiros e europeus. A escolha da rota Brasil–Europa fundamenta-se em três fatores principais. Em primeiro lugar, a União Europeia se destaca como o principal polo regulatório global no processo de descarbonização do transporte marítimo, com a implementação de instrumentos como o \gls{EUETS} e o regulamento FuelEU Maritime, que introduzem custos diretos associados às emissões de dióxido de carbono e metas progressivas de redução da intensidade de carbono dos combustíveis \citep{eu_reg_2023_957,eu_reg_2023_1805}.

Em segundo lugar, essa rota tem relevância econômica significativa, uma vez que a União Europeia figura entre os principais parceiros comerciais do Brasil, representando parcela expressiva do fluxo de comércio exterior e, consequentemente, do tráfego marítimo internacional. Por fim, a aplicação extraterritorial desses instrumentos regulatórios, que incidem sobre navios que escalam portos europeus independentemente do local de abastecimento, torna essa rota particularmente adequada para analisar, de forma integrada, os efeitos econômicos e ambientais associados à adoção de combustíveis de menor intensidade de carbono \citep{carvalho_biofuels_2021,eu_reg_2023_957}.

Desse modo, o recorte geográfico adotado não apenas reflete a relevância empírica da rota, mas amplia também a capacidade explicativa do modelo ao incorporar explicitamente sinais regulatórios que influenciam a viabilidade das alternativas energéticas no curto prazo.

No Brasil, a aplicação do biobunker carece de estudos integrados que avaliem seu desempenho sob premissas técnicas explícitas de compatibilidade operacional, articulando custos, impactos ambientais e condicionantes regulatórios com base em parâmetros operacionais representativos. Também não há consenso sobre sua inserção no setor a partir de instrumentos institucionais e regulatórios \citep{carvalho_biofuels_2021}. Essa lacuna orienta a presente pesquisa.

Diante dessas premissas, estabelece-se a seguinte pergunta de pesquisa: sob premissas técnicas de compatibilidade operacional, o biobunker pode ser considerado alternativa econômico-ambientalmente viável para apoiar a transição energética do transporte marítimo internacional operando a partir do Brasil, considerando as metas de descarbonização da \gls{IMO} até 2030 e os condicionantes regulatórios associados a rotas internacionais?

Para responder a essa pergunta, o objetivo geral é desenvolver um modelo determinístico de simulação comparativa \emph{in silico}, aplicado a uma rota representativa do transporte marítimo internacional operando a partir do Brasil, que permita avaliar a viabilidade econômico-ambiental do biobunker sob premissas técnicas de compatibilidade operacional. Essa avaliação considera os efeitos da fração de mistura e da rota de produção do biodiesel sobre o consumo requerido, as emissões no escopo \emph{well-to-wake}, o custo total por viagem e o custo marginal de abatimento de carbono, à luz das metas de descarbonização da \gls{IMO} até 2030 e dos condicionantes regulatórios associados à rota analisada.

Em termos argumentativos, a dissertação se organiza em torno de três contribuições principais. A primeira consiste em propor um arcabouço integrado para comparar cenários de biobunker a partir da articulação entre desempenho operacional agregado, emissões no escopo \emph{well-to-wake}, custo total por viagem e custo marginal de abatimento de carbono. A segunda consiste em aplicar esse arcabouço a uma rota representativa do transporte marítimo internacional operando a partir do Brasil, incorporando explicitamente a heterogeneidade entre rotas de produção do biodiesel. A terceira consiste em interpretar os resultados à luz do ambiente regulatório europeu incidente sobre rotas Brasil--Europa, distinguindo entre custo explícito de \emph{compliance} e parâmetro de referência para avaliação da atratividade econômico-ambiental.

Com base nessa estrutura, a dissertação trata a compatibilidade técnica do uso de misturas como premissa operacional delimitada no modelo e concentra sua verificação empírica na existência, ou não, de combinações entre fração de mistura e rota de produção do biodiesel que satisfaçam simultaneamente critérios ambientais e econômico-regulatórios. Essa operacionalização é formalizada no capítulo de método por meio de uma hipótese de existência de cenários viáveis no domínio analisado.

Além desta introdução, a dissertação está organizada em mais quatro capítulos. O Capítulo~\ref{chap_referencial} apresenta a revisão da literatura e o referencial teórico que sustentam a análise. O Capítulo~\ref{chap_metodo} descreve o método, as premissas de modelagem e os critérios de avaliação adotados. O Capítulo~\ref{chap_resultados} expõe os resultados parciais da simulação e sua discussão. Por fim, o Capítulo~\ref{chap_conclusao} reúne as considerações parciais da pesquisa e os encaminhamentos para as etapas seguintes do trabalho.
